\heading{高等数学A（II）-25春-期中}
\noteinfo{题目来源于赛艇，答案由AI生成。}

\begin{question}[计算积分]

\begin{enumerate}[label=(\arabic*),leftmargin=2.6em,itemsep=0.25em,topsep=0.2em]
\item 计算
\[
\iint_D\sqrt{|y-x^2|}\,dx\,dy,
\qquad D=[-1,1]\times[0,2].
\]
\item 计算
\[
\int_{-1}^{1}dx\int_{0}^{\sqrt{1-x^2}}dy\int_{1}^{1+\sqrt{1-x^2-y^2}}
\frac{dz}{\sqrt{x^2+y^2+z^2}}.
\]
\end{enumerate}

\begin{solution}

\begin{enumerate}[label=(\arabic*),leftmargin=2.6em,itemsep=0.25em,topsep=0.2em]
\item 先对 $y$ 积分，并按 $y=x^2$ 分段。计算结果为
\ansmath{\iint_D\sqrt{|y-x^2|}\,dx\,dy=\frac{5}{3}+\frac{\pi}{2}.}

\item 把区域改写成
\[
 x^2+y^2+(z-1)^2\le1,\qquad y\ge0,\qquad z\ge1.
\]
在以原点为极点的球坐标中，这相当于
\[
 0\le\theta\le\pi,
 \qquad 0\le\varphi\le\frac{\pi}{4},
 \qquad \sec\varphi\le\rho\le2\cos\varphi,
\]
其中 $\varphi$ 是与正 $z$ 轴的夹角。因被积函数为 $1/\rho$，故
\[
 I=\int_0^{\pi}\!d\theta\int_0^{\pi/4}\!\sin\varphi\,d\varphi
 \int_{\sec\varphi}^{2\cos\varphi}\rho\,d\rho.
\]
算得
\ansmath{I=\frac{\pi}{6}(7-4\sqrt2).}
\end{enumerate}

\end{solution}
\end{question}

\begin{question}[曲线积分]

计算曲线积分
\[
 I=\int_{\Gamma}\frac{(x-\tfrac12-y)\,dx+(x-\tfrac12+y)\,dy}{(x-\tfrac12)^2+y^2},
\]
其中 $\Gamma$ 从 $(0,-1)$ 沿抛物线 $y^2=1-x$ 到 $(0,1)$。

\begin{solution}

参数化取
\[
 x=1-t^2,\qquad y=t,\qquad -1\le t\le1.
\]
则
\[
 dx=-2t\,dt,\qquad dy=dt.
\]
代入后得到一个有理函数积分，直接计算可得
\ansmath{I=\frac{\pi}{2}+2\arctan 3.}

\end{solution}
\end{question}

\begin{question}[曲面积分]

计算曲面积分
\[
 I=\iint_{\Sigma}2xy\,dy\,dz-y^2\,dz\,dx-x^2\,dx\,dy,
\]
其中 $\Sigma$ 是抛物面 $x^2+y^2=z$ 在 $z=0$ 和 $z=1$ 之间的部分，取外侧。

\begin{solution}

把它看成向量场
\[
 \mathbf F=(2xy,-y^2,-x^2)
\]
穿过抛物面侧面的通量。其散度为
\[
 \nabla\cdot\mathbf F=\frac{\partial(2xy)}{\partial x}+\frac{\partial(-y^2)}{\partial y}+\frac{\partial(-x^2)}{\partial z}=2y-2y+0=0.
\]
把侧面与顶盖 $z=1$ 的圆盘补成闭曲面，由 Gauss 定理知：侧面通量等于顶盖通量的相反数。顶盖法向量向上，故
\[
 I_{\text{top}}=\iint_{x^2+y^2\le1}(-x^2)\,dxdy=-\frac{\pi}{4}.
\]
底部 $z=0$ 退化为一点，没有贡献，因此
\ansmath{I=-I_{\text{top}}=\frac{\pi}{4}.}

\end{solution}
\end{question}

\begin{question}[区域面积]

求曲线
\[
 x^3+y^3=xy
\]
所围成区域的面积。

\begin{solution}

该曲线关于直线 $y=x$ 对称，在第一象限形成一条闭曲线。令
\[
 y=tx,
\]
代入得
\[
 x^3(1+t^3)=tx^2\Longrightarrow x=\frac{t}{1+t^3},\qquad y=\frac{t^2}{1+t^3},\qquad t\in[0,\infty).
\]
利用参数曲线面积公式
\[
 S=\frac12\int (x\,dy-y\,dx),
\]
可得
\[
 S=\frac12\int_0^{\infty}\frac{t^2}{(1+t^3)^2}\,dt=\frac16.
\]
因此
\ansmath{S=\frac16.}

\end{solution}
\end{question}

\begin{question}[立体体积]

求旋转抛物面
\[
 z=6-x^2-y^2
\]
和圆锥面
\[
 z=\sqrt{x^2+y^2}
\]
所围成立体的体积。

\begin{solution}

改用柱坐标 $r=\sqrt{x^2+y^2}$。交线满足
\[
 6-r^2=r
 \Longrightarrow r=2.
\]
因此体积
\[
 V=\int_0^{2\pi}\!d\theta\int_0^2 (6-r^2-r)r\,dr
 =2\pi\int_0^2(6r-r^3-r^2)\,dr.
\]
计算得
\ansmath{V=\frac{32\pi}{3}.}

\end{solution}
\end{question}

\begin{question}[微分方程]

求以下方程的通解：
\begin{enumerate}[label=(\arabic*),leftmargin=2.6em,itemsep=0.25em,topsep=0.2em]
\item $y'-y=xy^5$;
\item $y''+y'=4\sin x+x\cos2x$。
\end{enumerate}

\begin{solution}

\begin{enumerate}[label=(\arabic*),leftmargin=2.6em,itemsep=0.25em,topsep=0.2em]
\item 对 $y\ne0$，令 $v=y^{-4}$，则
\[
 v'=-4y^{-5}y'=-4(v+x).
\]
即
\[
 v'+4v=-4x.
\]
解得
\[
 v=Ce^{-4x}-x+\frac14.
\]
故
\ansmath{y=(Ce^{-4x}-x+\tfrac14)^{-1/4}.}
另有平凡解 $y\equiv0$。

\item 齐次方程 $y''+y'=0$ 的通解为
\[
 y_h=C_1+C_2e^{-x}.
\]
取特解后可得
\ansmath{y=C_1+C_2e^{-x}+\frac{x\sin2x}{10}-\frac{x\cos2x}{5}-2\sin x-2\cos x+\frac{4}{25}\sin2x+\frac{13}{100}\cos2x.}
\end{enumerate}

\end{solution}
\end{question}

\begin{question}[Green 定理不等式]

设 $\Gamma$ 是取逆时针方向的圆周
\[
 (x-1)^2+(y-1)^2=1,
\]
且 $f(t)$ 为正的连续函数。证明：
\[
\oint_{\Gamma}x f(y)\,dy-\frac{y}{f(x)}\,dx\ge2\pi.
\]

\begin{solution}

令
\[
 P(x,y)=-\frac{y}{f(x)},\qquad Q(x,y)=xf(y).
\]
则由 Green 定理
\[
\oint_{\Gamma}P\,dx+Q\,dy
=\iint_D\left(\frac{\partial Q}{\partial x}-\frac{\partial P}{\partial y}\right)\,dA,
\]
其中 $D$ 是圆盘区域。这里
\[
 \frac{\partial Q}{\partial x}=f(y),
 \qquad
 -\frac{\partial P}{\partial y}=\frac{1}{f(x)}.
\]
所以
\[
\oint_{\Gamma}x f(y)\,dy-\frac{y}{f(x)}\,dx
=\iint_D\left(f(y)+\frac{1}{f(x)}\right)dA.
\]
由 $a+1/a\ge2$（对 $a>0$）得
\[
 f(y)+\frac{1}{f(x)}\ge2.
\]
而 $D$ 的面积为 $\pi$，故
\[
\oint_{\Gamma}x f(y)\,dy-\frac{y}{f(x)}\,dx
\ge2\iint_D1\,dA=2\pi.
\]
得证。

\ansmath{\oint_{\Gamma}x f(y)\,dy-\frac{y}{f(x)}\,dx
\ge2\iint_D1\,dA=2\pi}

\end{solution}
\end{question}

\begin{question}[单调性证明]

假设 $f(t)$ 为连续的正函数，证明：当 $t>0$ 时，
\[
F(t)=
\frac{\iiint_{x^2+y^2+z^2\le t^2}f(x^2+y^2+z^2)\,dxdydz}
{\iint_{x^2+y^2\le t^2}(x^2+y^2)f(x^2+y^2)\,dxdy}
\]
是严格单调递减函数。

\begin{solution}

把分子分母都改写成极坐标/球坐标积分。记
\[
 A(t)=\int_0^t r^2 f(r^2)\,dr,
 \qquad
 B(t)=\int_0^t r^3 f(r^2)\,dr.
\]
则
\[
\iiint_{x^2+y^2+z^2\le t^2}f(r^2)\,dV=4\pi A(t),
\qquad
\iint_{x^2+y^2\le t^2}(x^2+y^2)f(r^2)\,dA=2\pi B(t).
\]
因此
\[
 F(t)=2\frac{A(t)}{B(t)}.
\]
求导：
\[
 F'(t)=2\frac{A'(t)B(t)-A(t)B'(t)}{B(t)^2}
 =2\frac{t^2f(t^2)B(t)-t^3f(t^2)A(t)}{B(t)^2}.
\]
提取公因子得
\[
 F'(t)=\frac{2t^2f(t^2)}{B(t)^2}\Bigl(B(t)-tA(t)\Bigr).
\]
但对 $0\le r\le t$ 有 $r^3\le tr^2$，且 $f(r^2)>0$，故
\[
 B(t)=\int_0^t r^3f(r^2)\,dr
< t\int_0^t r^2f(r^2)\,dr=tA(t).
\]
从而
\[
 F'(t)<0.
\]
所以
\ansmath{F(t)\text{ 在 }t>0\text{ 上严格单调递减。}}

\end{solution}
\end{question}

